
\subsubsection{Interpretation: Folklore Validated with Perfect Evidence on MNIST}

Our results validate practitioner folklore for MNIST with quantitative rigor. Kaggle competition winners and PyTorch official tutorials implicitly avoid horizontal flip on MNIST. We now have clear answers for MNIST: horizontal flip degrades asymmetric digit accuracy by 0.72-1.00 percentage points at moderate flip rate (p=0.5), with dose-dependent degradation ranging from 0.37 pp (p=0.3) to 4.10 pp (p=0.9), while semantically valid augmentation (rotation $\pm 15$°) causes no differential effect. The perfect or near-perfect dose-response relationships (Spearman $\rho$=-1.0 or -0.969) indicate this is a deterministic causal mechanism.

The semantic validity framework demonstrated on MNIST provides practitioners with a principled criterion: an augmentation is valid if and only if the augmented image remains semantically in-distribution for the labeled class. Horizontal flip violates this criterion for asymmetric MNIST digits, introducing label noise that degrades test accuracy. Rotation $\pm 15$° satisfies the criterion, producing no degradation.

The digit 7 anomaly (minimal degradation -0.30\% versus digits 2/5 at -6.60\%/-6.93\%) suggests semantic validity operates on a continuum. Flipped digit 7 visually resembles canonical 7, reducing perceived label noise and degradation.

\subsubsection{Limitations}

We enumerate three principled limitations rooted in controlled experimental design:

\textbf{MNIST-Only Validation}: All experiments use MNIST handwritten digits ($28\times 28$ grayscale, 10 classes). Effect size and mechanism may differ on datasets with higher inter-class similarity, multi-channel color images, or different semantic constraints. We chose MNIST for rigorous proof-of-concept. The semantic validity principle is hypothesized to generalize, but effect sizes are dataset-dependent and require empirical validation.

\textbf{Standard CNN Architecture Only}: All experiments use a shallow convolutional network (~100K parameters). Modern architectures (ResNet, Vision Transformers, pre-trained models) may exhibit different robustness. Deeper models may learn more robust representations that mitigate label noise, or pre-trained models may have learned flip-invariance during pre-training. We deliberately chose a standard shallow CNN to test the effect in resource-constrained settings.

\textbf{Observational Design}: We infer the causal mechanism (flip $\rightarrow$ label noise $\rightarrow$ degradation) from dose-response correlation, not direct intervention. While Spearman $\rho$=-1.0 or -0.969 is exceptionally strong correlational evidence, gold-standard causal validation would require interventions such as flipping images but correcting labels. Nevertheless, our rotation control rules out alternative explanations, and the perfect dose-response provides correlational evidence stronger than most observational studies achieve.

\subsubsection{Broader Impact}

This work has positive societal impact by improving ML safety in critical domains. Medical imaging models trained with horizontal flip augmentation on lateralized pathology may silently degrade diagnostic accuracy. Traffic sign recognition systems trained with flip on directional arrows may misclassify critical signs. By demonstrating augmentation validation methodology on MNIST, we reduce deployment risk in safety-critical applications.

The semantic validity principle extends beyond augmentation to other modeling choices: data preprocessing, architecture design, and loss functions. Our dose-response validation methodology provides a template for responsible ML evaluation beyond the specific case of horizontal flip on MNIST.

